\begin{table*}[!t]
\centering
\footnotesize
\renewcommand{\arraystretch}{1.05}
\textbf{Nomenclature}\\[2pt]
\begin{tabular}{@{}L{0.2\textwidth}L{0.77\textwidth}@{}}
\toprule
\multicolumn{2}{@{}l}{\textbf{Decision and valuation variables}}\\
$I$ & Investment (retrofit) cost\\
$V^{*}$ & Investment trigger value\\
$\Vlow$, $\Vhigh$ & Trigger at the lower and upper endpoints of the ambiguity set $\Sigma_t$\\
$\Vlin$, $\VMEU$ & Member trigger under the affine rule and under the pre-commitment benchmark, at the same recorded weight\\
$\Vzero$ & Trigger at the volatility anchor $\sbar$\\
$V_0$ & Pre-commitment benchmark reference value\\
$V^{*}(\alpha,t)$ & Member-specific trigger under the affine rule, Eq.~(\ref{eq:affine})\\
$P^{*}$ & Marginal carbon-price threshold, Eq.~(\ref{eq:price})\\
$P^{*}_{i}$ & Marginal carbon-price threshold of member $i$, equal to $P^{*}$ evaluated at $\alpha_i$\\
$P^{*}_{i,\mathrm{lin}}$ & Marginal carbon-price threshold of member $i$ under the affine rule, used by the safeguard distances\\
$P^{*}_{0}$ & No-ambiguity marginal carbon-price threshold, called the anchor threshold\\
$P_t$ & Prevailing effective carbon price at gate $t$\\
$C_t$ & Contested-price interval, $[\min_i P^{*}_i,\max_i P^{*}_i)$. When all member thresholds coincide, this half-open interval is empty and $B_t=0$\\
$B_t$ & Width of the contested-price interval, Eq.~(\ref{eq:interval})\\
$\tau$ & First hitting time of a candidate trigger in the benchmark problem\\
$V_{\tau}$ & Project value at the first hitting time $\tau$\\
$J^{\alpha}$ & Pre-commitment $\aMEU$ benchmark objective, Eq.~(\ref{eq:Jalpha})\\
$c$, $s$, $\Delta$ & Common per-tonne accounting adjustments, with $\Delta=c-s$\\
$L_{\mathrm{lin}}$, $U_{\mathrm{lin}}$ & Lower and upper edges of the affine contested-price interval $C_t$\\[2pt]
\multicolumn{2}{@{}l}{\textbf{Stochastic process and economic parameters}}\\
$\sbar$ & Positive forward-volatility anchor, $\sbar>0$, estimated from history in the application\\
$\sigma_p$ & Carbon-price volatility, forward parameter of the GBM process\\
$\slow$, $\shigh$ & Lower and upper volatility endpoints of $\Sigma_t$\\
$\mu_c$ & Annual carbon-price growth scenario, declared as a decision input\\
$r$ & Discount rate\\
$\delta$ & Net discount margin, $\delta=r-\mu_c$\\
$Q$ & Annual captured emissions (t/yr), net of the capture rate\\
$\blo$, $\bhi$ & Positive roots of the characteristic equation at the low-volatility and high-volatility endpoints of the declared ambiguity set\\
$\beta_1$ & Positive root of the characteristic equation, Eq.~(\ref{eq:char}), as a function of $\sigma_p$\\
$\beta_0$ & Shorthand for $\beta_1(\sbar)$, the root at the volatility anchor\\
$\Db$ & Spread between the roots at the two volatility endpoints, $\Db:=\blo-\bhi\ge 0$, strictly positive at a positive radius (Proposition~\ref{prop:benchmark}). Distinct from the weight distance $d$ of the elicitation-precision target\\
$D_0$, $E_0$ & Auxiliary curvature terms in the closed-form $\tilde{\kappa}$, Eq.~(\ref{eq:kappa})\\[2pt]
\multicolumn{2}{@{}l}{\textbf{Ambiguity-radius and recorded-weight parameters}}\\
$\alpha$ & Recorded weight on the adverse endpoint under the $\alpha$-maxmin parameterization. It is elicited or configured before the gate and serves as a reporting input. At $\alpha=1$, the benchmark objective reduces to MEU\\
$\Sigma_t$ & Volatility ambiguity set, Eq.~(\ref{eq:ambset})\\
$h_t$ & Admissible ambiguity radius at gate $t$; in the application, $h_t=\min\{c_w w_t,(1-\epsilon)\sbar\}$\\
$c_w$ & Signal-to-radius scaling constant in the pre-cap map $h_t=c_w w_t$, $c_w\ge 0$\\
$\epsilon$ & Admissibility margin for the ambiguity-radius cap, $0<\epsilon<1$\\
$\rho$ & Relative ambiguity radius, $h/\sbar$. Marked variants denote other quantities, namely $\hat{\rho}_k$ (sample autocorrelation), Spearman $\rho$ (rank correlation), and $\rho_{\mathrm{retest}}$ (reliability)\\
$\tilde{\alpha}(h)$ & Critical weight at which the member trigger crosses the anchor trigger, for $h>0$, Eq.~(\ref{eq:critical}); $\tilde{\alpha}(0):=1/2$ by continuous extension\\
$\tilde{\alpha}_{\mathrm{lin}}$, $\tilde{\alpha}_{\mathrm{rec}}$, $\tilde{\alpha}_{\alpha\text{-MEU}}$ & Critical weight under the affine rule, the recursive rule, and the pre-commitment benchmark\\
$\sigma_{\mathrm{eff}}(\alpha)$ & Effective volatility of the recursive rule, $\sqrt{\alpha\slow^{2}+(1-\alpha)\shigh^{2}}$\\
$\alpha_{\min}$, $\alpha_{\max}$ & Smallest and largest recorded weights on a panel, which set the two interval edges\\
$\tilde{\kappa}(\sbar)$ & Volatility-channel curvature ratio, $V^{*\prime\prime}/(4V^{*\prime})$, Eq.~(\ref{eq:kappa})\\[2pt]
\multicolumn{2}{@{}l}{\textbf{Approximation and safeguard quantities}}\\
$\theta$ & Registered economic-calibration tuple $(r,\mu_c,\sbar)$\\
$\eR(\theta,\rho)$ & Registered refined-grid maximum of the normalized error between affine and benchmark member triggers, Eq.~(\ref{eq:safeguard})\\
$\ER(\theta,\rho)$ & Price-unit release envelope specific to a calibration and a radius, Eq.~(\ref{eq:safeguard})\\
$\dedge$, $\dmember$ & Price-edge and nearest-member distances used by the safeguard, Eq.~(\ref{eq:safeguard})\\
$\varepsilon_{\mathrm{num}}$ & Numerical tolerance added to the registered maximum error\\[2pt]
\multicolumn{2}{@{}l}{\textbf{Policy-text input quantities}}\\
$w_t$ & Quarterly rolling variability signal, the four-quarter population standard deviation of PSDI, Eq.~(\ref{eq:wt})\\
$\PSDI_t$ & Policy-semantic-dispersion index, Eq.~(\ref{eq:psdi})\\
$\bar{c}_t$ & Centroid of normalized document embeddings in period $t$\\[2pt]
\multicolumn{2}{@{}l}{\textbf{Abbreviations}}\\
$\aMEU$ & Alpha-maxmin expected utility\\
PSDI & Policy-semantic-dispersion index\\
\bottomrule
\end{tabular}
\par\smallskip
\parbox{0.97\textwidth}{\emph{Notes.} Symbols appearing in at least one numbered equation or at least six times in the mathematics are listed. The Supplementary Material contains the complete list of 93 entries.}
\end{table*}


\section{Literature review}

A gate report built from member thresholds and their contested-price interval relates to three research streams. Ambiguity-aware real options provide the trigger analysis, committee decision support provides the reporting context, and text-based uncertainty measures provide one application input. Each stream supplies a different object to the gate and leaves one gap for member-level screening.

\subsection{Real-options analytics for irreversible investment}

\citet{McDonaldSiegel1986} demonstrated the positive value of waiting under irreversibility and uncertainty. \citet{DixitPindyck1994} developed this result into the trigger analysis used here. \citet{NadarajahSecomandi2023} surveyed the uptake of this framework in the operations literature on energy assets. In this literature, the valuation target is an operational asset, and the flexibility is managerial. The retrofit examined here is such an asset, and the quantity of interest is the date of commitment. In the standard perpetual geometric Brownian motion (GBM) model, this date depends on volatility, because higher volatility raises the value of waiting and the trigger. When volatility is ambiguous in the sense of \citet{Ellsberg1961}, more than one law for the price process remains admissible. The trigger then depends on how this ambiguity is represented and evaluated.

\citet{NishimuraOzaki2007} placed Knightian ambiguity on the drift of the price process in an irreversible investment problem. They evaluated it under the maxmin expected utility (MEU) preferences of \citet{GilboaSchmeidler1989}. In their model, the worst-case prior reduces the value of waiting and accelerates commitment. \citet{TrojanowskaKort2010} extended the drift-ambiguity analysis to a finite project life, where acceleration held only in the short run. The $\alpha$-maxmin generalization \citep{Ghirardato2004} assigns weight $\alpha$ to the worst-case prior and $1-\alpha$ to the best-case prior, nesting MEU at $\alpha=1$. Under this broader formulation, \citet{Schroder2011} showed how the investment decision depends on both the ambiguity and the attitude toward it. Depending on this attitude, the subjective project value can rise, and commitment can become more attractive. Two decision makers facing the same objective situation may therefore act differently. \citet{Schroder2020} noted a further issue: ambiguity models of this class, including $\aMEU$, are generally not dynamically consistent, a property needed for recursive optimal stopping. The direction of the effect of ambiguity on exercise timing also varies with the setting. In \citet{MiaoWang2011}, ambiguity could accelerate or delay option exercise, depending on whether it affected continuation payoffs, termination payoffs, or both. Drift and volatility enter the characteristic equation through different terms, so the comparative statics for drift ambiguity require rederivation under volatility ambiguity. Given the dynamic inconsistency of $\aMEU$, the benchmark in this paper is a pre-commitment rule. It is compared with a recursive alternative, which is time-consistent by construction.

Other formulations place ambiguity on volatility or on the model more broadly, and the decision question of each determines the object it returns. Four objects recur, and one formulation may return more than one of them. Formulations concerned with the worth of a claim return a value. The uncertain-volatility model of \citet{Avellaneda1995} bounds the price of a claim by evaluating it at the extremes of a volatility band. The $G$-expectation calculus of \citet{Peng2019} gives this band a coherent dynamic form, and \citet{EpsteinJi2013} carried the construction into asset pricing. Formulations concerned with how an agent should act return a policy. \citet{Matoussi2015} solved robust utility maximization in nondominated models through second-order backward stochastic differential equations. \citet{Kallblad2018} obtained dynamically consistent investment criteria under model uncertainty through robust forward preferences. Formulations concerned with stopping, such as multiple-prior and time-consistent treatments, return a stopping region or an equilibrium policy \citep{Riedel2009,HuangYu2021}. Formulations concerned with commitment return a single robust decision. Recursive multiple priors belong to this group \citep{ChenEpstein2002,EpsteinSchneider2003}. So do robust control \citep{HansenSargent2001} and distributionally robust optimization \citep{MohajerinKuhn2018}. They evaluate uncertainty at the level of the decision problem and return one decision robust over the whole set. A robust recursion returns the same type of object, and a rectangularity condition preserves its Bellman structure \citep{Iyengar2005}. The ambiguity set may also be a confidence region estimated from observed data \citep{Wiesemann2013}, and the reported output remains one worst-case policy. Under smooth ambiguity \citep{Klibanoff2005}, \citet{MazzonTankov2024} reduced a coal-plant divestment problem with learning to a family of standard stopping problems. Their output is still the policy of a single decision maker. Values, policies, regions, and robust decisions are all indexed by the problem, so none of the reviewed formulations returns the price range of member disagreement.

One strand of the ambiguity literature works with sets of valuations and comes closer to such a range. In the Knightian decision theory of \citet{Bewley2002}, incomplete preferences represent ambiguity. One act ranks above another only if its expected utility is higher under every prior in the set. Some pairs therefore remain unranked, and an inertia assumption resolves them in favor of the status quo. An alternative is then adopted only when it dominates the status quo unanimously. \citet{RigottiShannon2005} carried this structure into general equilibrium, where incomplete preferences induce a set of valuations. The unanimity principle of \citet{Bewley2002} suggests an analogous reporting logic for stage-gate screening, because deferral is the incumbent option at a gate. A unanimous invest classification here records agreement among the recorded member rules at the prevailing price. Where unanimity fails, aggregation is one possible downstream step. Axiomatic treatments of this step exist for heterogeneous beliefs and ambiguity attitudes. \citet{Gilboa2004} restricted the Pareto principle to choices under identical beliefs. \citet{Cres2011} aggregated maxmin experts sharing a utility function. \citet{Danan2016} applied the Pareto criterion only to comparisons robust to belief imprecision. Each of these studies characterized a social criterion, whereas the member-threshold vector precedes the aggregation step.

The contested-price interval is computed from the member-threshold vector, whose entries are indexed by recorded weights. A project-level formulation can supply such a vector only after an explicit member-indexing step. For the pre-commitment objective, this step solves the same scalar problem at each recorded weight. A time-consistent formulation yields scalar thresholds only if the equilibrium stopping region of every member has a unique boundary in the same ordered state variable. The ambiguity set used here meets this condition.

The applied literature on CCUS retrofit timing also works at the project level. \citet{Compernolle2017} evaluated offshore enhanced oil recovery with CO$_2$ storage. \citet{Fan2023} and \citet{Li2026} examined the retrofit timing of coal-fired plants, and \citet{Sun2024} examined incentive design. \citet{Colombe2024} worked at the system level and derived a deployment plan through two-stage stochastic programming. \citet{Deeney2021} built a real-options decision support tool for research and development investment in CO$_2$ recycling. Their compound option structure carries technical and market risk and values the project.

By contrast, one line of theory indexes the option by the individual. \citet{Garlappi2017,Garlappi2022} examined group real options in which members hold heterogeneous beliefs and decide by voting. In their model, a representative median member cannot reproduce group behavior. Their analysis concerns group behavior after a vote, whereas the present report precedes any vote. The present setting differs from this literature in three further respects. Heterogeneity enters through ambiguity attitudes, all members share the same information set, and ambiguity falls on volatility. \citet{Cardenas2025} extended this line to the timing of group investment decisions. In this model, members hold heterogeneous beliefs, and the rule of each member has a double continuation region. Disagreement can produce inertia and underinvestment even when every member would invest individually. The two boundaries delimit the continuation region of one member, whereas the interval reported here spans the one-sided thresholds of different members. The gate reports this interval without resolving the disagreement. Across this stream, the reported objects are project-level triggers, values, policies, stopping regions, robust decisions, and group behavior after a vote. In the reviewed work, member-level rules appear only inside group models of voting or joint timing \citep{Garlappi2017,Cardenas2025}. The first gap is therefore the absence of a gate-level vector of one-sided member thresholds indexed by ambiguity attitude. This gap includes the position of such thresholds relative to the no-ambiguity anchor under volatility ambiguity.

\subsection{Stage-gate screening and committee-level decision support}

Stage-gate processes organize investment evaluation as sequential committee reviews \citep{Cooper2008}. The supporting decision systems typically deliver aggregate scores, rankings, or recommendations \citep{Sprague1980,Power2002}, which are updated as conditions change \citep{Violos2025}. Group decision methods such as TOPSIS-based scoring \citep{HwangYoon1981} also compress heterogeneous evaluations into a single score or ranking. A score ranks alternatives at one decision point, whereas a member-specific threshold is the state-variable level at which the screening classification of a member changes. A committee can use both, with the thresholds locating the prevailing price relative to each member and a scoring procedure ranking the remaining alternatives. The interval width measures dispersion among model-implied screening thresholds in price units. By contrast, consensus measures, reviewed by \citet{HerreraViedma2014}, assess alignment among stated preferences before a collective choice.

A separate operational-research tradition treats disagreement as an output of the analysis. \citet{Phillips1984} built the decision model jointly with the problem owners and counted it as requisite once further sensitivity analysis generated no new intuitions. \citet{FrancoMontibeller2010} developed this stance into facilitated modeling for groups. The present method offers this tradition a quantitative device for one setting, translating recorded weights into thresholds without a further aggregation step. Two groups of recent methods operate downstream of the resulting report: aggregation procedures and procedures deciding above the individual member. \citet{Dong2026} integrated information across attributes and periods to produce superiority-probability comparisons and probabilistic rankings. \citet{SongLi2026} aggregated numerical and linguistic preference relations at the subgroup level and then solved minimum-adjustment consensus models under limited compromise. Portfolio and multi-stage models kept the decision above the individual member \citep{Marques2022,ShenLi2023,Saiz2024}. Recent robust and adaptive designs did so as well \citep{Shen2025,ShavazipourStewart2026}. They use an uncertainty set to select one policy with acceptable performance across the set. These studies address what the organization should commit to, whereas the present method examines whether the members screening one project currently agree. For each member, it reports the signed margin, defined as the member threshold minus the prevailing price. The margin is positive under a wait classification. At the committee level, the method reports the interval edges. Aggregating the edges away at the gate would leave a single number without a record of the member behind it. A later consensus model, ranking, or portfolio selection could use the member thresholds, given a stated rule for how they enter its objectives and constraints. Most reviewed methods aggregate or reconcile member assessments before reporting, and facilitated modeling lacks a quantitative threshold device. No reviewed method converts recorded member weights into thresholds on the price axis, and this missing report is the second gap.

\subsection{Text-based uncertainty measures as a model input}

In the application, only one gate input, the ambiguity radius, comes from policy text. Existing studies commonly derive text-based uncertainty measures from keyword frequencies or supervised classifiers. The Economic Policy Uncertainty index counts prespecified newspaper terms \citep{Baker2016}. Investment research often uses it as a regression covariate \citep{GulenIon2016}. \citet{Hassan2019} extended this design to firm-level political risk, and \citet{Gavriilidis2021} applied it to climate news. Fixed term lists may miss paraphrases and emerging terminology. \citet{Hartley2025} reported closer agreement with human-audited labels for large language model classifiers than for dictionary rules. Whether policy documents carry a usable dispersion signal remained open in this comparison. Supervised alternatives, such as the MacBERT-based climate policy uncertainty index for China of \citet{Ma2023}, require labeled training data. The limits of fixed term lists and labeled data motivate the unsupervised, embedding-based construction in the application. This construction follows the text-as-data approach reviewed by \citet{Gentzkow2019} and yields a series measuring semantic dispersion across contemporaneous policy documents. In the reviewed studies, text-based measures served as indices, document labels, or regression covariates. None of them supplied an input to a member-level screening rule, which is the third gap.

Table~\ref{tab:lit} summarizes the three gaps. First, the reviewed real-options formulations return project-level objects or group behavior, and none reports a gate-level vector of member thresholds indexed by ambiguity attitude. Second, committee decision support aggregates member assessments into scores, rankings, or consensus measures, so the price range of member disagreement stays unreported. Third, text-based uncertainty measures have not been linked to member-level screening rules. The present method addresses the first two gaps with a member-indexed threshold vector on the price axis. The critical weight locates this vector relative to the anchor, and a stated accuracy standard sets its release conditions. For the third gap, the application supplies a replaceable text-based radius as one gate input.

\begin{table*}[!t]
\caption{Relation of prior literature streams to the contested-price interval and its width. The first two rows are the streams underlying the method, and the third row covers the application-specific input of the running example.}
\label{tab:lit}
\centering
\footnotesize
\renewcommand{\arraystretch}{1.15}
\begin{tabular}{@{}L{0.27\textwidth}L{0.18\textwidth}L{0.22\textwidth}L{0.25\textwidth}@{}}
\toprule
Literature stream & Contribution & Gap for stage-gate screening & Role in this paper\\
\midrule
Ambiguity-aware real options \citep{NishimuraOzaki2007,TrojanowskaKort2010,Schroder2011,Fan2023} & Investment timing rules under non-unique priors & Gate-level vector of member thresholds indexed by ambiguity attitude & Analytical basis for member triggers\\
Decision support and stage-gate screening \citep{Sprague1980,Cooper2008,ShenLi2023,Saiz2024,Shen2025,Violos2025} & Sequential review processes and aggregate scores or rankings & Explicit reporting of member-threshold dispersion from recorded weights & Interface logic reporting the contested-price interval, its width, and the current classification summary\\
Application input. Text-based policy-uncertainty measurement \citep{Baker2016,Gavriilidis2021,Ma2023,Hartley2025} & Quantitative signals from policy and news text & Integration of such signals into member-specific gate-level screening rules & Construction of the replaceable input $w_t$\\
\bottomrule
\end{tabular}
\end{table*}
